\documentclass[11pt]{article}

\usepackage[T1]{fontenc}
\usepackage{lmodern}
\usepackage[a4paper,margin=28mm]{geometry}
\usepackage{amsmath,amssymb,amsthm,mathtools,mathrsfs}
\usepackage{microtype}
\usepackage{enumitem}
\usepackage{booktabs,array}
\usepackage{xcolor}
\usepackage[colorlinks=true,linkcolor=blue!48!black,
  citecolor=blue!48!black,urlcolor=blue!55!black]{hyperref}
\setlength{\emergencystretch}{2em}

\newtheorem{theorem}{Theorem}[section]
\newtheorem{proposition}[theorem]{Proposition}
\newtheorem{corollary}[theorem]{Corollary}
\newtheorem{lemma}[theorem]{Lemma}
\theoremstyle{definition}
\newtheorem{definition}[theorem]{Definition}
\theoremstyle{remark}
\newtheorem{remark}[theorem]{Remark}

\newcommand{\C}{\mathbb C}
\newcommand{\R}{\mathbb R}
\newcommand{\N}{\mathbb N}
\newcommand{\cA}{\mathcal A}
\newcommand{\cI}{\mathcal I}
\newcommand{\cS}{\mathcal S}
\newcommand{\cT}{\mathcal T}
\newcommand{\cQ}{\mathcal Q}
\newcommand{\cW}{\mathcal W}
\newcommand{\cR}{\mathcal R}
\newcommand{\cH}{\mathcal H}
\newcommand{\Tr}{\operatorname{Tr}}
\newcommand{\im}{\operatorname{im}}
\newcommand{\Span}{\operatorname{span}}
\newcommand{\Rea}{\operatorname{Re}}
\newcommand{\ord}{\operatorname{ord}}
\newcommand{\id}{\operatorname{id}}
\newcommand{\mul}{\operatorname{mul}}

\title{\textbf{One orientation survives}\\
\large An orientation--transgression completeness law for the Riemann source}
\author{Michael Haufschild}
\date{}

\begin{document}

\maketitle

\begin{abstract}
We construct and classify a two-step arithmetic source built from the
Meyer--Calder\'on defect and the directed algebraic rational-Witt packet.
At each prime \(p\), degree one is
\[
  \cA_p^1=\cQ_p\oplus\cW_p,
\]
degree two is the first-suspension packet \(\cA_p^2\), and the
Calder\'on--Witt cross-effect gives a vector-valued alternating cup
\[
  \Lambda^2\cA_p^1\longrightarrow\cA_p^2.
\]
The cup is onto, its vector-valued radical is zero, and its kernel
\(\cI_p\) is the exact incidence relation space.  We give an explicit
section and Green retraction, proving
\[
  \Lambda^2\cA_p^1
    =\cI_p\oplus\im(\cA_p^2),\qquad
  H^\bullet\!\left(
    \cI_p\to\Lambda^2\cA_p^1\to\cA_p^2\right)=0.
\]
The same assertions hold on the all-prime finite-support source.

The arithmetic source nevertheless has one canonical absolute invariant.
A normalized Haar coefficient \(\tau:\cA^2\to\C\) defines the oriented
boundary and trace
\[
  D w=(\smile w,-\tau(\smile w)),\qquad
  T(z,c)=\tau(z)+c.
\]
Then \(\im D=\ker T\), the oriented source has cohomology
\[
                         (H^0,H^1,H^2)=(0,0,\C),
\]
and every top cochain has a canonical Green decomposition
\[
                         z=DHz+\jmath Tz.
\]
This yields the paper's universal law.  For every complex vector space
\(V\) and linear consumer \(F\),
\[
 F D=0
 \quad\Longleftrightarrow\quad
 \exists!\,\phi:\C\to V\ \text{such that}\ F=\phi T.
\]
Thus every regular boundary-blind consumer sees only one orientation
line.  Information beyond it is exactly sensitivity to an exact boundary.
After restricting incidence primitives to a sector \(\cS\), the complete
remaining invariant is the canonical Green primitive modulo \(\cS\);
its class vanishes exactly when the Green primitive remains in \(\cS\).

The law unifies previously separate completion, scalarization,
hypertranscendence, Rees, and positivity obstructions.  It also changes the
Riemann frontier: another source, scalar trace, or regular Hilbert
completion cannot supply the missing information.  The only surviving
mechanism is a source-pointed singular polarized chain transport which
preserves the vector-valued cup through the boundary and makes its Green
primitive sector-preserving.  In the completed divisor geometry this
condition is the vanishing of the Frobenius--Hodge commutator; its
weight-one metric law is equivalent to the Riemann hypothesis.  No
zero-exclusion statement is assumed or claimed.
\end{abstract}

\medskip
\noindent\textbf{Keywords.}
Riemann zeta function; arithmetic source; rational Witt vectors;
Calder\'on transform; relative cohomology; Green contraction;
orientation; Hodge polarization; singular completion.

\medskip
\noindent\textbf{2020 Mathematics Subject Classification.}
Primary 11M26, 11R58; Secondary 18G35, 46A13, 47A06.

\section{Introduction}

The persistent ambiguity in spectral approaches to the Riemann hypothesis
is not whether the Riemann divisor can be represented.  It can.  The
difficulty is to retain the signed displacement
\[
                         \Rea\rho-\frac12
\tag{1.1}\label{eq:signed-displacement}
\]
while passing from an arithmetic source to a completed positive or
spectral object.

Several exact mechanisms show where the sign is lost.  A scalar
codiagonal preserves total mass but annihilates the labelled zero-mass
radical.  Product-topology graph closure retains every finite output and
becomes maximally vertical.  Reflection-even squaring preserves a divisor
but erases its normal orientation.  A positive Ward index counts an
off-line pair without forcing the count to vanish.  Reflected
prime-valuation sectors can have incompatible asymptotic types, so no
finite-order graded map sews them.  These statements appear different
because they are expressed in Banach, Hilbert, divisor, probabilistic, and
Rees coordinates.

The present paper identifies their common source.  The starting object is
not a Hilbert space selected after the divisor is known.  It is the
directed algebraic union of finite rational-Witt packets, coupled to the
one-dimensional Meyer--Calder\'on defect.  The coupling is bilinear and
therefore belongs as multiplication in a two-step graded source.  Its
alternating polarization is vector-valued:
\[
  (q,a)\smile(q',a')
     =\chi(q,a')-\chi(q',a).
\tag{1.2}\label{eq:cup-intro}
\]
Keeping this value before scalarization is decisive.

The source turns out to be much more rigid than expected.  Its exterior cup
generates all of degree two.  The kernel is an exact incidence space, and
an explicit Green operator contracts the resulting three-term complex.
Thus no hidden absolute middle class remains to be discovered.  Adding the
normalized Haar orientation does not create an uncontrolled new
cohomology: it leaves exactly one top line and nothing else.

This produces one universal classification:
\[
\boxed{
\begin{gathered}
\text{regular absolute information}
   =\text{one source orientation line},\\[2mm]
\text{all further information}
   =\text{canonical primitive escaping an admitted sector}.
\end{gathered}}
\tag{1.3}\label{eq:master-law-intro}
\]
The first line is a factorization theorem for every vector-valued linear
consumer.  The second line is a necessity-and-sufficiency theorem for every
relative cocycle.  Neither is a dimension count or a heuristic.

\subsection*{What the law answers}

The law closes the source-identification problem.  It proves that one
should not search for another absolute source class, another scalar
positive functional, or another regular completion.  Those categories are
already classified.  A construction that carries more than the orientation
line must detect an exact boundary, hence must be relative to a restricted
primitive sector.

For the Riemann problem the relevant restrictions are simultaneously
arithmetic, topological, asymptotic, and polarizing.  The primitive must
retain prime, valuation depth, rational-Witt stage, suspension, and
source-horizon sample labels.  Its completion must be finer than every
finite-coordinate observation and more singular than every fixed Rees
order.  Its degree-two value must remain vector-valued until Hodge
polarization.  These are not optional design preferences: each is forced
by an exact obstruction.

The completed divisor geometry supplies the final coordinate.  If
\(\Theta\) is the weight operator, \(J\) the odd Hodge involution,
\(\smile\) the oriented divisor cup, and
\[
                       g(x,y)=x\smile Jy,
\tag{1.4}\label{eq:hodge-metric-intro}
\]
then
\[
 g(\Theta x,y)+g(x,\Theta y)-g(x,y)
   =-x\smile[\Theta,J]y.
\tag{1.5}\label{eq:purity-defect-intro}
\]
The cup is nondegenerate, so exact weight-one conservation is equivalent
to \([\Theta,J]=0\).  For the completed Riemann divisor, this is equivalent
to the Riemann hypothesis.  The source law therefore identifies the
missing mathematical object without presenting an RH equivalence as its
novelty: it is a singular, sector-preserving, polarized chain transport
from the explicit arithmetic source to this divisor metric.

\subsection*{Organization}

Section~\ref{sec:algebraic-source} constructs the algebraic source.
Section~\ref{sec:milnor} gives its exterior presentation and Green
contraction.  Section~\ref{sec:orientation} adds the unique orientation
line.  Section~\ref{sec:main-law} proves the universal
orientation--transgression law.  Sections~\ref{sec:shadows}
and~\ref{sec:riemann-boundary} derive the earlier completion,
scalarization, Rees, and Hodge results as shadows of the law.

\section{The directed algebraic source}
\label{sec:algebraic-source}

\subsection{Finite rational-Witt packets}

Fix a prime \(p\).  Let \(\mathfrak S_p\) be the directed set of finite
Frobenius stages.  At a stage \(s\in\mathfrak S_p\), let \(X_{p,s}\) be its
finite affine packet and
\[
                         W_{p,s}=\C^{X_{p,s}}.
\tag{2.1}\label{eq:finite-stage}
\]
Gauge-corrected refinement gives injective linear pullbacks
\[
                  \kappa_{s,t}^{*}:W_{p,s}\longrightarrow W_{p,t}
                  \qquad(s\le t)
\tag{2.2}\label{eq:gauge-pullback}
\]
which are reflexive and strictly composable.  All stages embed in one
full-support Haar \(L^2\)-envelope.  Define the algebraic packet
\[
                  \cW_p=\bigcup_{s\in\mathfrak S_p}W_{p,s}.
\tag{2.3}\label{eq:algebraic-packet}
\]
Directedness has a useful exact consequence: every \(a\in\cW_p\) occurs at
one finite stage.  No closure-generated vector is part of the arithmetic
source.

\begin{proposition}[Finite-stage realization]
\label{prop:finite-stage-realization}
Every \(a\in\cW_p\) has a stage \(s\) and a unique stage image
\(a_s\in W_{p,s}\) representing it in the Haar envelope.  Refinement of
\(a_s\) represents the same element at every later stage.
\end{proposition}

\begin{proof}
The stage images form a directed family of subspaces.  Membership in their
supremum is therefore membership in one member of the family.  Injectivity
of the stage maps gives uniqueness at a fixed stage, and strict
composability of \(\kappa_{s,t}^{*}\) gives refinement invariance.
\end{proof}

\subsection{The two degree-one legs}

Let \(\cQ_p\) be the Meyer--Calder\'on quotient.  It is one-dimensional
and has a distinguished unit defect \(u_p\).  Let
\[
                        \lambda_p:\cQ_p\overset{\sim}{\longrightarrow}\C
\tag{2.4}\label{eq:calderon-coordinate}
\]
be the normalized coordinate with \(\lambda_p(u_p)=1\).

Let
\[
                         S_p:\cW_p\longrightarrow \cA_p^2
\tag{2.5}\label{eq:first-suspension}
\]
be the first-suspension map, with codomain restricted to its algebraic
image.  This map is a linear isomorphism onto \(\cA_p^2\).  Let
\(\delta_p\in\R_{>0}\) be the nonzero polarized half-density scale.

Define
\[
   \cA_p^1=\cQ_p\oplus\cW_p,\qquad
   \chi_p(q,a)=\lambda_p(q)\,\delta_p\,S_p(a),
\tag{2.6}\label{eq:two-step-source}
\]
and
\[
  (q,a)\smile_p(q',a')
    =\chi_p(q,a')-\chi_p(q',a).
\tag{2.7}\label{eq:local-cup}
\]
This is bilinear, alternating, and invariant under every
gauge-corrected stage refinement.

\begin{theorem}[Vector-valued source nondegeneracy]
\label{thm:vector-source-nondegenerate}
The left and right radicals of \(\smile_p\) are zero.  The all-prime
finite-support cup
\[
  \smile:\cA^1\times\cA^1\longrightarrow\cA^2,\qquad
  \cA^j=\bigoplus_p\cA_p^j,
\tag{2.8}\label{eq:global-cup}
\]
also has zero vector-valued radical.
\end{theorem}

\begin{proof}
Suppose \(x=(q,a)\) cups to zero with every \(y\).
Taking \(y=(0,b)\) gives
\(\lambda_p(q)\delta_pS_p(b)=0\) for every \(b\).
The packet is nonzero, \(S_p\) is injective, and \(\delta_p\ne0\), so
\(q=0\).  Taking \(y=(u_p,0)\) then gives
\(-\delta_pS_p(a)=0\), hence \(a=0\).  Skew-symmetry gives the right
radical.  In the all-prime direct sum, testing against vectors supported at
one prime reduces each component to the local argument.
\end{proof}

\section{Exterior presentation and Green contraction}
\label{sec:milnor}

Alternation induces a linear exterior cup
\[
                 d_1:\Lambda^2\cA_p^1\longrightarrow\cA_p^2.
\tag{3.1}\label{eq:exterior-cup}
\]
Put
\[
                           \cI_p=\ker d_1.
\tag{3.2}\label{eq:incidence}
\]
We call \(\cI_p\) the Milnor-incidence relation space.  The name records
its role: it contains exactly the quadratic degree-one incidences which
vanish in the suspension center.

\subsection{Surjectivity and a canonical section}

For \(z\in\cA_p^2\), define
\[
 a_z=\delta_p^{-1}S_p^{-1}z,\qquad
 \sigma_p(z)=u_p\wedge(0,a_z).
\tag{3.3}\label{eq:cup-section}
\]
Then
\[
                              d_1\sigma_p=\id_{\cA_p^2}.
\tag{3.4}\label{eq:section-law}
\]
Consequently \(d_1\) is onto and
\[
             \cA_p^2\cong\Lambda^2\cA_p^1/\cI_p.
\tag{3.5}\label{eq:milnor-quotient}
\]

Define the Green retraction
\[
      h_p(w)=w-\sigma_p(d_1w)\in\cI_p.
\tag{3.6}\label{eq:green-retraction}
\]
Writing \(\iota_p:\cI_p\hookrightarrow\Lambda^2\cA_p^1\) for inclusion,
\[
                   \iota_ph_p+\sigma_pd_1=\id.
\tag{3.7}\label{eq:green-middle}
\]

\begin{theorem}[Exact Milnor resolution]
\label{thm:milnor-resolution}
For every prime \(p\),
\[
  \Lambda^2\cA_p^1
     =\cI_p\oplus\im\sigma_p,
\tag{3.8}\label{eq:milnor-splitting}
\]
and the three-term complex
\[
   0\longrightarrow\cI_p
     \xrightarrow{\ \iota_p\ }
     \Lambda^2\cA_p^1
     \xrightarrow{\ d_1\ }
     \cA_p^2
     \longrightarrow0
\tag{3.9}\label{eq:milnor-complex}
\]
is split exact.  Its cohomology vanishes in all three degrees.

For the all-prime finite-support source, the analogous exterior cup is
onto, cross-prime wedges lie in its incidence kernel, an explicit
finite-support section satisfies the same Green identity, and the global
three-term complex is split exact.
\end{theorem}

\begin{proof}
Equation~\eqref{eq:section-law} gives surjectivity.
Equation~\eqref{eq:green-retraction} gives \(h_p(w)\in\ker d_1\), and
\eqref{eq:green-middle} gives the sum decomposition.  If
\(w\in\cI_p\cap\im\sigma_p\), write \(w=\sigma_pz\).  Then
\(0=d_1w=z\), so \(w=0\).  Hence the sum is direct.
Injectivity of \(\iota_p\), exactness at the middle term by definition of
\(\cI_p\), and surjectivity of \(d_1\) kill \(H^0,H^1,H^2\).

For the direct sum, every degree-two vector has finite prime support.  Sum
the local sections over that support.  Local cups recover each component,
while cross-prime wedges have zero cup and therefore are incidence
relations.  The same calculation proves the global Green identity.
\end{proof}

\begin{remark}
\label{rem:not-hidden-cohomology}
The theorem rules out a specific class of explanations.  The Riemann
divisor cannot be an undiscovered absolute class of the algebraic
Calder\'on--Witt Milnor resolution: that resolution is contractible.
Any additional class must be created by an orientation extension, a
restriction of primitives, or a singular completion.
\end{remark}

\section{The unique absolute orientation}
\label{sec:orientation}

\subsection{Haar trace}

Let \(1_p\in\cW_p\) be the constant packet and put
\[
  m_p=\langle1_p,1_p\rangle,\qquad
  \eta_p(a)=m_p^{-1}\langle1_p,a\rangle.
\tag{4.1}\label{eq:packet-trace}
\]
Full support gives \(m_p\ne0\), and \(\eta_p(1_p)=1\).  Transport through
\(S_p\) defines
\[
                     \tau_p:\cA_p^2\longrightarrow\C.
\tag{4.2}\label{eq:center-trace}
\]
The constant packet gives a section
\(\epsilon_p:\C\to\cA_p^2\) with
\[
                           \tau_p\epsilon_p=\id_\C.
\tag{4.3}\label{eq:trace-section}
\]

\subsection{Oriented graph complex}

Define
\[
\begin{aligned}
 D_p &: \Lambda^2\cA_p^1\longrightarrow\cA_p^2\oplus\C,
 &D_p w&=(d_1w,-\tau_p(d_1w)),\\
 T_p &: \cA_p^2\oplus\C\longrightarrow\C,
 &T_p(z,c)&=\tau_p(z)+c.
\end{aligned}
\tag{4.4}\label{eq:oriented-graph}
\]
Then \(T_pD_p=0\).

\begin{theorem}[Oriented source cohomology]
\label{thm:oriented-cohomology}
The graph boundary and orientation trace form an exact sequence
\[
 \Lambda^2\cA_p^1\xrightarrow{D_p}
 \cA_p^2\oplus\C\xrightarrow{T_p}\C\longrightarrow0,
\qquad
 \im D_p=\ker T_p.
\tag{4.5}\label{eq:oriented-exact}
\]
Moreover,
\[
 \ker D_p=\cI_p,\qquad
 H^\bullet\!\left(
  \cI_p\to\Lambda^2\cA_p^1\xrightarrow{D_p}
  \cA_p^2\oplus\C\right)=(0,0,\C).
\tag{4.6}\label{eq:oriented-cohomology}
\]
The all-prime finite-support oriented source has the same cohomology.
\end{theorem}

\begin{proof}
If \(T_p(z,c)=0\), surjectivity of \(d_1\) gives \(w\) with \(d_1w=z\).
The trace equation gives \(c=-\tau_p(z)\), hence \((z,c)=D_pw\).
The reverse inclusion follows from \(T_pD_p=0\).

The first component of \(D_pw\) is \(d_1w\), so
\(\ker D_p=\ker d_1=\cI_p\).  Exactness and injectivity at the first two
terms give \(H^0=H^1=0\).  The cokernel of \(D_p\) is canonically \(\C\)
through \(T_p\), giving \(H^2\cong\C\).

The global proof uses the finite-support sum
\(\tau=\sum_p\tau_p\).  Every prime separately supplies a section of this
trace, so it is onto.  Global cup surjectivity gives
\(\im D=\ker T\), and the remaining assertions follow as above.
\end{proof}

Define the orientation inclusion
\[
                     \jmath_p(c)=(0,c)
\tag{4.7}\label{eq:orientation-inclusion}
\]
and the top Green lift
\[
                     H_p(z,c)=\sigma_p(z).
\tag{4.8}\label{eq:top-green-lift}
\]

\begin{proposition}[Top Green decomposition]
\label{prop:top-green}
Every top cochain has the canonical decomposition
\[
                   D_pH_p(z,c)+\jmath_pT_p(z,c)=(z,c).
\tag{4.9}\label{eq:top-green}
\]
The same identity holds globally.
\end{proposition}

\begin{proof}
Using \(d_1\sigma_pz=z\), the first coordinate is \(z\).  The second
coordinate is
\[
               -\tau_p(z)+\tau_p(z)+c=c.
\]
The global calculation is identical.
\end{proof}

\section{The orientation--transgression completeness law}
\label{sec:main-law}

We first isolate the universal algebra behind
\eqref{eq:top-green}.  Let
\[
 A_1\xrightarrow{D}A_2\xrightarrow{T}O
\tag{5.1}\label{eq:abstract-complex}
\]
be linear maps over a field.  Suppose there are linear maps
\(\jmath:O\to A_2\) and \(H:A_2\to A_1\) such that
\[
                  DH+\jmath T=\id_{A_2},\qquad
                  TD=0,\qquad T\jmath=\id_O.
\tag{5.2}\label{eq:abstract-green}
\]

\begin{lemma}[Universal orientation factorization]
\label{lem:universal-factorization}
For every vector space \(V\) and linear map \(F:A_2\to V\),
\[
  FD=0
  \quad\Longleftrightarrow\quad
  \exists!\,\phi:O\to V\text{ such that }F=\phi T.
\tag{5.3}\label{eq:factorization}
\]
Equivalently,
\[
  F\text{ does not factor through }T
  \quad\Longleftrightarrow\quad
  \exists x\in A_1,\ F(Dx)\ne0.
\tag{5.4}\label{eq:boundary-detection}
\]
\end{lemma}

\begin{proof}
Assume \(FD=0\).  Applying \(F\) to \eqref{eq:abstract-green} gives
\[
                         F=(F\jmath)T.
\]
Thus \(\phi=F\jmath\) is a factor.  If \(F=\psi T\), then
\[
                         \psi=\psi T\jmath=F\jmath=\phi,
\]
so the factor is unique.

Conversely, \(F=\phi T\) gives \(FD=\phi TD=0\).  Negating the two
equivalent statements yields \eqref{eq:boundary-detection}.
\end{proof}

The middle-degree counterpart is relative.  Let
\[
 A_0\xrightarrow{d_0}A_1\xrightarrow{d_1}A_2
\tag{5.5}\label{eq:three-term}
\]
have a contraction \(h_1:A_1\to A_0\), \(h_2:A_2\to A_1\) satisfying
\[
                         d_0h_1+h_2d_1=\id_{A_1}.
\tag{5.6}\label{eq:middle-contraction}
\]
For an admitted primitive sector \(\cS\subseteq A_0\), define
\[
 \mathfrak o_{\cS}(x)
   =[h_1x]\in A_0/(\cS+\ker d_0).
\tag{5.7}\label{eq:relative-obstruction}
\]

\begin{lemma}[Relative primitive classification]
\label{lem:relative-classification}
If \(d_1x=0\), then
\[
 \mathfrak o_{\cS}(x)=0
 \quad\Longleftrightarrow\quad
 \exists b\in\cS,\ d_0b=x.
\tag{5.8}\label{eq:relative-existence}
\]
If \(d_0\) is injective, this is equivalent to
\[
                  \mathfrak o_{\cS}(x)=0
                  \quad\Longleftrightarrow\quad
                  h_1x\in\cS.
\tag{5.9}\label{eq:sector-preservation}
\]
\end{lemma}

\begin{proof}
Because \(d_1x=0\), \eqref{eq:middle-contraction} gives
\(d_0h_1x=x\).  If \([h_1x]=0\), write
\(h_1x=b+k\) with \(b\in\cS\) and \(k\in\ker d_0\).  Then
\(d_0b=x\).  Conversely, if \(b\in\cS\) and \(d_0b=x\), then
\(h_1x-b\in\ker d_0\), so \([h_1x]=0\).

If \(d_0\) is injective, its kernel vanishes.  The quotient class is zero
exactly when \(h_1x\in\cS\).
\end{proof}

\begin{theorem}[Source orientation--transgression completeness]
\label{thm:main}
For the local and all-prime algebraic Calder\'on--Witt sources, the
following classification holds.

\begin{enumerate}[label=\textup{(\roman*)},leftmargin=2.2em]
\item The un-oriented Milnor source is split exact.  The oriented source has
      cohomology \((0,0,\C)\), with canonical generator measured by its
      normalized Haar trace \(T\).

\item For every complex vector space \(V\) and every linear top-source
      consumer \(F\),
      \[
      FD=0
      \quad\Longleftrightarrow\quad
      \exists!\,\phi:\C\to V,\ F=\phi T.
      \tag{5.10}\label{eq:arithmetic-factorization}
      \]
      Hence information beyond the orientation trace is equivalent to
      detecting an exact oriented boundary.

\item For every admitted incidence sector
      \(\cS\subseteq\cI\) and every exterior cocycle \(x\),
      \[
      \mathfrak o_{\cS}(x)=0
      \quad\Longleftrightarrow\quad
      hx\in\cS.
      \tag{5.11}\label{eq:arithmetic-relative}
      \]
      Thus the relative class is exactly canonical Green-primitive escape.
\end{enumerate}

Consequently the source has no third regular information channel:
all absolute boundary-invariant information is the one orientation line,
and all additional information is relative transgression.
\end{theorem}

\begin{proof}
Part~(i) is Theorems~\ref{thm:milnor-resolution}
and~\ref{thm:oriented-cohomology}.  Proposition~\ref{prop:top-green}
supplies \eqref{eq:abstract-green}; the oriented trace kills the graph
boundary and the orientation inclusion is its section.  Applying
Lemma~\ref{lem:universal-factorization} proves part~(ii).

For part~(iii), use \(d_0=\iota:\cI\hookrightarrow\Lambda^2\cA^1\),
\(d_1=\smile\), \(h_1=h\), and \(h_2=\sigma\).
Equation~\eqref{eq:green-middle} is the contraction identity and
\(\iota\) is injective.  Lemma~\ref{lem:relative-classification} gives
\eqref{eq:arithmetic-relative}.
\end{proof}

\begin{corollary}[No hidden regular divisor channel]
\label{cor:no-third-channel}
A linear source observable which is insensitive to exact incidence
boundaries contains no divisor information independent of the normalized
orientation trace.  Any proposed additional source observable must either
detect an exact boundary or depend on a restriction of the primitive
sector.
\end{corollary}

\begin{proof}
This is Theorem~\ref{thm:main}(ii)--(iii).
\end{proof}

\section{Earlier obstructions as shadows of one law}
\label{sec:shadows}

This section does not assemble independent main theorems.  It identifies
which operation in Theorem~\ref{thm:main} is violated by each known
obstruction.

\subsection{Scalarization before polarization}

The local degree-one source has a one-dimensional Calder\'on leg and a
packet leg containing two linearly independent continuous directions.
Let
\[
                    \ell:\cA_p^2\longrightarrow\C
\tag{6.1}\label{eq:scalar-trace}
\]
be any scalar trace, and define
\[
 \omega_\ell((q,a),(q',a'))
     =\ell\!\left((q,a)\smile_p(q',a')\right).
\tag{6.2}\label{eq:scalar-cup}
\]

\begin{proposition}[Scalar-trace fracture]
\label{prop:scalar-fracture}
For every prime \(p\) and every scalar trace \(\ell\), the alternating form
\(\omega_\ell\) is degenerate.
\end{proposition}

\begin{proof}
Fix the Calder\'on unit \(u_p\).  If \(\omega_\ell\) were nondegenerate,
the map
\[
  a\longmapsto \ell\chi_p(u_p,a)
\tag{6.3}\label{eq:packet-scalar-map}
\]
would be injective on the packet leg.  But the packet leg contains a
two-dimensional subspace and the target is one-dimensional.  This is
impossible.
\end{proof}

Thus the unique top orientation trace of
Theorem~\ref{thm:main} must not be confused with a scalarization of each
degree-one leg.  Multiplication must occur while the degree-two packet is
still vector-valued; only afterward may an orientation be read.

\subsection{Closed completion and maximal verticality}

Let source representatives \(x_K\to0\) carry output packets whose leading
finite coordinate projections are all onto.  If a closed linear relation
\(\cR\) in product topology contains all pairs \((x_K,y_j)\), then
closedness gives \((0,y_j)\in\cR\).  Surjectivity of every leading
projection makes \(\mul\cR\) dense; closedness makes it the entire target.
Hence
\[
                       \cR=\operatorname{dom}\cR\times Y.
\tag{6.4}\label{eq:maximal-vertical}
\]
The completed relation transports no output information.

In the language of Theorem~\ref{thm:main}, product closure has admitted so
many boundary directions that the Green splitting is erased.  A useful
completion must use a topology strictly finer than finite-coordinate
product observation.

\subsection{Finite-order reflection sewing}

Let \(a_N\to0\).  Suppose one reflected source germ is all-orders flat,
\[
                      x_N=o(a_N^m)\quad\text{for every }m,
\tag{6.5}\label{eq:flat-germ}
\]
while the other has finite grade
\[
                      \|y_N\|\asymp a_N^d.
\tag{6.6}\label{eq:finite-grade}
\]
Every exact sewing \(R_Nx_N=y_N\) satisfies
\[
                  a_N^{-K}=o(\|R_N\|)
                  \quad\text{for every }K.
\tag{6.7}\label{eq:beyond-rees}
\]
Thus the required relative transgression is beyond every fixed Rees order.
The reflected prime-valuation source realizes precisely this split at an
off-critical displacement.

\subsection{Hypertranscendental coefficient sectors}

A translation cocycle can be exact in the ambient meromorphic field while
having no primitive in a differentially algebraic coefficient sector.
The cocycle is not a new ambient class.  It is the class of its canonical
ambient primitive modulo the thin sector.  This is
\eqref{eq:relative-obstruction} literally, with differential algebraicity
as the admitted primitive condition.

\subsection{Positive singular examples}

Two source constructions show that boundary sensitivity need not be
pathological.

First, splitting the Euler common-test source at the moving point
\(\log n\) gives a compact Gamma/co-Poisson head and a Hardy tail.  In the
joint scaling limit their kernels converge to complementary shifted-polar
channels whose sum reconstructs the Hardy Gram exactly.  The polar channel
is created by boundary blow-up; it is not an external correction.

Second, multiplication of the Euler boundary half-density by the
source-derived factor \(\sqrt{\log p}\) produces the critical
von-Mangoldt/Mellin current before scalar fusion.  Its reflection-odd
center jet is the exterior difference of squared source heights.  Gamma
anchoring centers those squares, and every finite anchored packet is
nonnegative.  This realizes the minimal source promotion forced by the
first-clock capacity bound, while its cofinal control remains necessarily
singular.

\section{The Riemann boundary coordinate}
\label{sec:riemann-boundary}

\subsection{Derived polarization}

Let \(V\) be a complex vector space with a nondegenerate alternating cup
\(\omega\), a complex-linear weight operator \(\Theta\), and a
conjugate-linear odd polarization \(J\).  Assume
\[
 \omega(\Theta x,z)+\omega(x,\Theta z)=\omega(x,z).
\tag{7.1}\label{eq:cup-weight}
\]
Define
\[
                            g(x,y)=\omega(x,Jy).
\tag{7.2}\label{eq:derived-metric}
\]

\begin{proposition}[Purity-defect identity]
\label{prop:purity-defect}
For all \(x,y\in V\),
\[
 g(\Theta x,y)+g(x,\Theta y)-g(x,y)
    =-\omega\bigl(x,[\Theta,J]y\bigr).
\tag{7.3}\label{eq:purity-defect}
\]
Consequently,
\[
\begin{aligned}
 &g(\Theta x,y)+g(x,\Theta y)=g(x,y)
       \quad\text{for all }x,y\\
 &\hspace{40mm}\Longleftrightarrow\quad [\Theta,J]=0.
\end{aligned}
\tag{7.4}\label{eq:metric-commutator}
\]
\end{proposition}

\begin{proof}
Apply \eqref{eq:cup-weight} to \(x\) and \(Jy\):
\[
 \omega(\Theta x,Jy)+\omega(x,\Theta Jy)=\omega(x,Jy).
\]
Subtract and add \(\omega(x,J\Theta y)\).  This gives
\eqref{eq:purity-defect}.  If the commutator vanishes, the metric law
follows.  Conversely, the metric law and nondegeneracy of \(\omega\) force
\([\Theta,J]y=0\) for every \(y\).
\end{proof}

\subsection{Completed divisor specialization}

For the completed Riemann divisor, the oriented cup is nondegenerate and
the Hodge involution pairs a divisor point with its reflected point.
Orbitwise,
\[
 [\Theta,J]_{\rho}
   =-\kappa(\rho)\,J_{\rho},\qquad
 \kappa(\rho)=2\Rea\rho-1.
\tag{7.5}\label{eq:orbit-curvature}
\]
The same coefficient has three source representations:
\[
\begin{aligned}
\kappa(\rho)
 &= -\bigl(h_p(\rho)-h_p(1-\overline\rho)\bigr),\\
\kappa(\rho)
 &= -\,\text{reflected prime-valuation entropy at }\rho,\\
-m_\rho(\rho-\tfrac12)
 &=\operatorname*{Res}_{s=\rho}
     \left[-(s-\tfrac12)\frac{\Xi'(s)}{\Xi(s)}\right].
\end{aligned}
\tag{7.6}\label{eq:charge-representations}
\]
The first identity is independent of the prime base, the second records
the Rees type, and the third retains the full complex outer charge.

\begin{theorem}[Riemann Hodge boundary law]
\label{thm:riemann-hodge}
For the completed Riemann divisor metric,
\[
 \begin{aligned}
 &g(\Theta x,y)+g(x,\Theta y)=g(x,y)
      \quad\text{for all }x,y\\
 &\hspace{38mm}\Longleftrightarrow
      [\Theta,J]=0
      \Longleftrightarrow
      \Rea\rho=\frac12
      \text{ for every nontrivial zero }\rho.
 \end{aligned}
\tag{7.7}\label{eq:rh-hodge}
\]
\end{theorem}

\begin{proof}
The first equivalence is
Proposition~\ref{prop:purity-defect}.  By
\eqref{eq:orbit-curvature}, the commutator vanishes on the orbit of
\(\rho\) exactly when \(2\Rea\rho-1=0\).  It vanishes globally exactly
when this holds for every nontrivial zero.
\end{proof}

\subsection{What is now missing}

Theorems~\ref{thm:main} and~\ref{thm:riemann-hodge} locate the boundary
between source construction and zero exclusion.

The arithmetic source already supplies:

\begin{enumerate}[leftmargin=2.2em]
\item a directed finite-stage packet before completion;
\item a nondegenerate vector-valued exterior product;
\item its exact incidence relations and Green contraction;
\item one normalized absolute orientation line;
\item a primitive promotion to the critical von-Mangoldt current; and
\item a singular Gamma/Hardy boundary decomposition.
\end{enumerate}

The divisor side already supplies:

\begin{enumerate}[leftmargin=2.2em]
\item the unsquared outer charge;
\item the odd Hodge polarization;
\item the exact purity-curvature identity;
\item the prime-valuation asymptotic type of that curvature; and
\item an RH-equivalent weight-one conservation law.
\end{enumerate}

The missing object is therefore not another item on either list.  It is one
chain-level transport between them.  The classification forces it to have
the following form:
\[
\boxed{
\begin{gathered}
\text{rational-Witt exterior source}\\
\xrightarrow[\text{nonlocal, beyond all finite Rees orders}]
 {\text{singular source-pointed Green transport}}\\
\text{Gamma/co-Poisson reflected boundary}\\
\xrightarrow[\text{after vector-valued multiplication}]
 {\text{odd Hodge polarization}}
\text{weight-one divisor metric}.
\end{gathered}}
\tag{7.8}\label{eq:missing-transport}
\]

It must preserve the canonical Green primitive in the admitted
positive/Hodge sector.  By
\eqref{eq:arithmetic-relative}, that is exactly vanishing of the relative
source class.  By \eqref{eq:rh-hodge}, its divisor realization is exactly
the vanishing of \([\Theta,J]\).  This is a specification of one
mathematical object, not a request for a new positivity criterion.

\section{Consequences for proof architecture}

\subsection{Necessary order of operations}

The source law fixes the order in which operations may occur:
\[
\begin{gathered}
\text{finite rational-Witt refinement}\\
\longrightarrow
\text{labelled degree-one source}\\
\longrightarrow
\text{vector-valued exterior multiplication}\\
\longrightarrow
\text{singular reflected boundary transport}\\
\longrightarrow
\text{Hodge polarization}\\
\longrightarrow
\text{orientation trace}.
\end{gathered}
\tag{8.1}\label{eq:operation-order}
\]
Reversing any adjacent pair meets an established obstruction.
Scalarizing before multiplication gives
Proposition~\ref{prop:scalar-fracture}.  Closing before installing an
unbounded sample grade gives \eqref{eq:maximal-vertical}.  Sewing reflected
sectors with finite order violates \eqref{eq:beyond-rees}.  Squaring before
odd polarization erases \eqref{eq:signed-displacement}.

\subsection{A classification of failed consumers}

The four principal consumer classes can now be separated without examining
zeros.

\begin{center}
\begin{tabular}{>{\raggedright\arraybackslash}p{0.24\linewidth}
                >{\raggedright\arraybackslash}p{0.31\linewidth}
                >{\raggedright\arraybackslash}p{0.34\linewidth}}
\toprule
Consumer & Preserved source datum & Forced loss\\
\midrule
bounded label-blind codiagonal
  & total orientation mass
  & labelled zero-mass radical\\
closed product-topology relation
  & every finite output packet
  & all output constraints\\
scalar trace before cup
  & one scalar degree-two moment
  & nondegeneracy of degree one\\
finite-Rees reflected sewing
  & finite asymptotic grades
  & all-orders-flat reflected sector\\
\bottomrule
\end{tabular}
\end{center}

These are not four missing tricks.  They are four ways of forcing a
boundary-sensitive observable through the regular factorization
\eqref{eq:arithmetic-factorization}.

\subsection{Positive consumers}

The law also explains why several positive constructions remain useful.
The Euler common-test source retains labels before the horizon split.  The
Gamma/co-Poisson head is a genuine positive Gram channel.  The Hardy tail
is a second positive channel.  Their joint blow-up restores the shifted
polar component.  The primitive-promoted exterior source keeps both legs
until an odd center jet is taken.  These constructions respect the order
\eqref{eq:operation-order}; their limitation is not source absence but
cofinal sector preservation.

\section{Discussion}

\subsection{A source, not a decoration}

The algebraic source is independent of any chosen zero.  Its finite-stage
packets, cup, incidence relations, trace, and Green contraction are defined
before the Xi divisor enters.  The one-dimensional top cohomology is
therefore not a decorative rewriting of the functional equation.  It is
the complete regular quotient of an arithmetic multiplication law.

\subsection{Why one orientation line is enough}

At first sight, one scalar top class seems too small to control a complex
divisor.  The factorization theorem says the opposite.  The scalar is not
expected to encode every divisor coordinate directly.  It classifies what
survives after exact boundaries have been forgotten.  Divisor displacement
belongs to the relative part: it is carried by whether the canonical
primitive stays inside the admitted sector.  The absolute line normalizes
orientation; relative transgression supplies location.

\subsection{Universality}

The factorization argument uses only a Green decomposition, a trace which
kills boundaries, and a section of that trace.  It is valid over arbitrary
fields and with arbitrary target dimension.  The relative classification
uses only a three-term contraction.  Arithmetic enters by producing these
structures canonically from rational-Witt and Calder\'on data.

The resulting universality is explanatory rather than terminological.
Source-horizon verticality, dense Hilbert radicals, scalar-trace rank loss,
hypertranscendental anomalies, Rees superselection, and Hodge purity are
corollaries or representations of the same dichotomy:
\[
                \text{orientation factor}\quad\text{or}\quad
                \text{primitive escape}.
\tag{9.1}\label{eq:dichotomy}
\]

\section{Conclusion}

The arithmetic source problem has a complete answer.  The directed
rational-Witt packet and Meyer--Calder\'on defect form a two-step graded
source with nondegenerate vector-valued cup.  Its exterior square has an
exact incidence presentation and explicit Green contraction.  Without
orientation it is contractible.  With normalized Haar orientation it has
exactly one top class.

Every regular boundary-blind consumer factors uniquely through that class.
Every additional invariant is the canonical primitive modulo an admitted
sector.  This is the orientation--transgression completeness law.

For the Riemann problem, the law removes a broad ambiguity.  Missing
mathematics is not another source, divisor equivalence, positive scalar
form, or regular completion.  It is one singular source-pointed polarized
chain transport which preserves the Green primitive through the
Gamma/co-Poisson boundary.  The completed divisor metric identifies its
decisive condition: weight-one conservation, equivalently vanishing of the
Frobenius--Hodge commutator.  The law does not prove that condition.  It
does prove why all known successful pieces and failed regular routes are
shadows of it, and what a source-internal proof must construct.

\end{document}
